
% To derive the dynamical model of the vehicle, we start by modeling the
% force and torque vectors caused by a single propeller in the body
% frame denoted~$\mathcal{B}$, whose origin is coincident with the
% vehicle's center of mass (CoM).

% In this section we model the propulsion system, starting with the
% force and torque contributions to the vehicle body of single
% propeller, and then summing up the combining the contributions of them
% all.

Consider a single propeller $i$ whose motor is rigidly linked to the
body frame~$\mathcal{B}$, as depicted in Figure~\ref{fig:prop}.  This
propeller  generates a reaction force (thrust) $\bar{F}_i$ and  torque
$\bar{M}_i$ on the vehicle body.
While the former results directly from the
propeller thrust,
\begin{equation}
  \label{eq:Fi}
  \bar{F}_i = f_i \hat{u}_i \qquad %\quad\text{for}\quad
  f_i = K_1 u_i
\end{equation}
where $u_i$ is the actuation signal and $f_i$ is the scalar thrust,
the later is the sum of two components: the torque caused by the
non-central thrust
and the propeller reaction torque:
\begin{equation}
  \label{eq:Mi}
  \bar{M}_i = \bar{r}_i\times\bar{F}_i - \tau_i \hat{u}_i \qquad
  \tau_i = w_i K_2 u_i
\end{equation}
where $\times$ denotes vector cross product and $w_i$ is $-1$ or $1$
depending on whether the propeller rotates clockwise or anti-clockwise
for a positive forward thrust $f_i>0$.  These linear relationships
result from the momentum-blade element theory by considering $u=n^2$,
where $n$ is the blade's rotation speed (in revolutions per second) while the
constants $K_1$ and $K_2$ given by
\begin{equation}
  \label{eq:blade}
  K_1 = \rho D^4 C_T \qquad K_2 = \frac{\rho D^5}{2\pi} C_P
\end{equation}
where $\rho$ is the air density, $D$ is the propeller diameter, and
the $C_T$ and $C_P$ are blade dependent adimensional constants called
thrust and power coefficients~\cite{mccormick95}.

The relative position $\bar{r}_i$ of the $i$-th propeller, orthogonal to the
$Z$ axis, and the unit vector $\hat{u}_i$, aligned with the propeller
axis, are uniquely defined by the angles $\theta_i$ and $\phi_i$, and can be
easily obtained from geometrical reasoning by:
\begin{equation}
  \bar{r}_i = \left(
    \begin{array}{c}
      d\cos(\theta_i) \\
      d\sin(\theta_i) \\
      0
      \end{array}
    \right)
    \qquad
    \hat{u}_i = \left(
      \begin{array}{c}
        \sin(\theta_i)\sin(\phi_i) \\
        -\cos(\theta_i)\sin(\phi_i) \\
        \cos(\phi_i)
      \end{array}
      \right)
\end{equation}
where $d=\|\bar{r}_i\|$ is the distance from the propeller to the CoM.
Stacking together the resulting force $\bar{F}_i$ and torque
$\bar{M}_i$, one can obtain the linear relation
\begin{equation}
  \label{eq:fm}
  \left(
    \begin{array}{c}
      \bar{F}_i \\
      \bar{M}_i
    \end{array}
  \right) = \bar{a}_i\,u_i
\end{equation}
where
\begin{equation}
  \begin{split}
    \bar{a}_i  &= \left(
      \begin{array}{c}
        K_1 \hat{u}_i \\
        K_1 \bar{r}_i \times \hat{u}_i - w_i K_2 \hat{u}_i
      \end{array}
    \right) \\
    &= \left(
      \begin{array}{c}
        K_1 \sin(\theta_i)\sin(\phi_i) \\
        -K_1 \cos(\theta_i)\sin(\phi_i) \\
        K_1 \cos(\phi_i) \\
        \left[K_1 d \cos(\phi_i) - w_i K_2 \sin(\phi_i)\right]\sin(\theta_i) \\
        -\left[K_1 d \cos(\phi_i) - w_i K_2 \sin(\phi_i)\right]\cos(\theta_i) \\
        -K_1 d \sin(\phi_i) - w_i K_2 \cos(\phi_i)
      \end{array}
    \right)
    \end{split}
\end{equation}

\begin{figure}
  \centering
  \includegraphics[width=0.6\linewidth]{prop.pdf}
  \caption{Notation used for modeling a single propeller with respect
    to the body frame~$\mathcal{B}$, centered on the vehicle's CoM.}
  \label{fig:prop}
\end{figure}

For a vehicle with $N$ propellers, the resulting force and torque will
be given by the sum of the contributions of each propeller, given with
respect to the CoM, by~\eqref{eq:fm}. This sum can be put in matrix
form as
\begin{equation}
  \label{eq:au}
  \left(
    \begin{array}{c}
      \bar{F} \\
      \bar{M}
    \end{array}
  \right) = \mathbf{A}\,\bar{u}
\end{equation}
where $\mathbf{A}=[\bar{a}_1 \cdots \bar{a}_N]$ is a square matrix,
hereby called \emph{actuation matrix}, and $\bar{u}=[u_1\cdots u_N]^T$
is the actuation input vector. \emph{The crucial observation is that, if the
actuation matrix $\mathbf{A}$ has at least rank 6, the linear
equation~\eqref{eq:au} can be solved for $\bar{u}$ for any given
combination of $\bar{F}$ and $\bar{M}$.}  A necessary (but not
sufficient\footnote{Sufficiency required $\mathbf{A}$ to be full
  rank.}) condition for this to be true is to have at least 6
propellers, thus justifying the hexarotor design considered in this
paper. Note that this matrix only depends on the vehicle's design
parameters, defined by the angles $\{\theta_i\}$ and $\{\phi_i\}$,
distance $d$, the trust coefficients $K_1$, $K_2$, and $\{w_i\}$.


%%% Local Variables:
%%% mode: latex
%%% TeX-master: "main"
%%% End:
